\documentclass{article}
\usepackage{graphicx} % Required for inserting images
\usepackage[a4paper,margin=1in]{geometry}
\usepackage[T1]{fontenc}
\usepackage[utf8]{inputenc}
\usepackage{hyperref}
\usepackage{enumitem}
\usepackage{lmodern}
\usepackage{amsmath}
\usepackage{comment}
\usepackage{xcolor}
\usepackage{float}
\usepackage{microtype}

\setlist[itemize]{leftmargin=1.5em}
\setlist[enumerate]{leftmargin=1.5em}

\newcommand{\notimpl}[1]{\textcolor{red}{#1}}
\title{Authority Document \\Urban Capability Model}
\date{\today}

\begin{document}
\maketitle

\section{Introduction}

\subsubsection*{What is this document?}
This document is the authority document for the Urban Capability Model. It details the objects that make up the model, the relationships that connect those objects, and the code that implements them. It is the single, standalone reference for the model's structure and behavior.

\subsubsection*{What does the model do?}
The model measures how easily people living in a city can reach the things that support everyday life and wellbeing, given the transport options available to them.

A \emph{point of interest} (\emph{POI}) is a real place that offers something people need: a supermarket, a pharmacy, a park, a school canteen. Each \emph{POI} belongs to a \emph{POI type}, the category it falls under (for example ``supermarket'' or ``pharmacy''). The model lets an analyst define which \emph{POI types} to track, group \emph{POI types} into \emph{services}, and group \emph{services} into \emph{capabilities}. Our case study groups its services into three capabilities: restorativeness (e.g. nature contact), nutrition (e.g. food access), and care (e.g. primary care).

The model also lets an analyst choose which transport modes to evaluate: walking, cycling, driving, and public transport. For many points spread across the study area, the model works out how easily each one can reach the relevant POIs by the enabled modes, and combines those into a single accessibility figure per \emph{POI type}. It then aggregates \emph{POI type} figures into a \emph{service} score, and \emph{service} scores into a \emph{capability} rating using a method called ELECTRE TRI, which sorts each capability into one of five ordinal ratings.

The output is a set of capability ratings for every evaluated point in the city. These ratings feed into an interactive map, where users can explore how each capability varies across the area and inspect which POIs and services drove the rating at any given location.

\subsubsection*{What is the purpose of this document?}
This document describes the model at a level that does not require reading the source code, while remaining precise enough to serve as a reference for anyone extending, auditing, or reusing the model. It separates what the model contains (Section 2, Object content), how those contents interact (Section 3, Dynamics), and the implementation choices behind each step (Section 4, Implementation), so a reader can consult only the level of detail they need.

\subsubsection*{How is the document structured?}
The document has three main sections. Object content introduces the objects that make up the model: the inputs, configuration, intermediate results, and outputs, and what each one represents. Dynamics explains how those objects interact, describing how the model turns inputs into results step by step. Implementation details the implementation choices behind each step.

\section{Object content}
The object model is organised around a three-level domain hierarchy: \textbf{POI types} are aggregated into \textbf{services}, and services are aggregated into \textbf{capabilities}. The following section denotes the structure of each object in the system : what it represents and what different fields make up the state of that object. In the document, we distinguish objects from fields based on the style of the font. For example, this is an \emph{object} and this is a \texttt{field}.
\subsection{Input objects}
\subsubsection*{Boundary of extraction}
    The \emph{boundary} defines the geographic area the model analyzes. It has two fields: a \texttt{multipolygon\_\allowbreak geometry}, the shape of the area, and a \texttt{coordinate\_\allowbreak reference\_\allowbreak system}, the standard that gives meaning to the geometry's coordinates (for example WGS84). The model dissolves the \texttt{multipolygon\_\allowbreak geometry} into one single shape. This shape sets the extraction area for the transport networks and the origin points of the analysis.

    The model allows two ways to set the study area: an analyst can provide a \emph{boundary} directly, or provide a \texttt{city\_name} in the \emph{general configuration}, in which case the model derives the boundary and extracts the networks automatically from OpenStreetMap data.
\subsubsection*{POI input collection}
    The \emph{POI input} is a collection of point, line, and polygon feature sets that
    share one common schema. Each \emph{POI} record carries four fields:
    \begin{itemize}
        \item \texttt{poi\_id}, which uniquely identifies the \emph{POI} record.
        \item \texttt{geometry}, the source geometry of the \emph{POI}.
        \item \texttt{coordinate\_reference\_system}, the spatial reference system
        of the source geometry.
        \item the classification attribute identified by
        \texttt{poi\_label\_field} in the \emph{general configuration}, whose value is
        used to associate the record with one or more configured
        \emph{POI types}.
    \end{itemize}
    A point geometry gives the model one candidate access location. A line or polygon geometry gives the model several candidate access locations, one for each vertex of the line or of the polygon's boundary.

    The model allows two ways to obtain \emph{POI} data: an analyst can supply the point, line, and polygon feature sets directly, or provide a \texttt{city\_name} in the \emph{general configuration}, in which case the model downloads the \emph{POI} data automatically from OpenStreetMap.

\subsubsection*{General Transit Feed Specification (GTFS)}
\emph{GTFS}, the General Transit Feed Specification, is a standard format for public-transport schedules and their geographic information. Its principal entities are:
\begin{itemize}
    \item \texttt{agency}: the public-transport operator included in the feed.
    \item \texttt{stops}: the stops and stations where passengers board or alight.
    \item \texttt{routes}: the public-transport lines the operator runs.
    \item \texttt{trips}: the individual journeys performed along each route.
    \item \texttt{stop\_times}: the sequence and timing of stops for each trip.
    \item \texttt{calendar}: when each service runs.
    \item \texttt{shapes}: the geographic path a trip follows.
\end{itemize}
The model accepts more than one \emph{GTFS} feed at once. An analyst could, for example, combine a feed describing the subway network with a separate feed describing the bus network, so a single analysis accounts for both.

\subsection{Configuration objects}
\subsubsection*{General configuration}
The \emph{general configuration} identifies one analysis run and binds its internal parameters. Its fields are grouped as follows:
\begin{itemize}
    \item \emph{study area identity}:
    \begin{itemize}
        \item \texttt{city\_name}, the name of the city the model uses to derive the study area when it does not read the boundary from a shapefile.
        \item \texttt{use\_shapefile}, a flag that tells the model which of the two ways to set the study area to use. When true, the model reads the study area boundary from the shapefile named by \texttt{name\_\allowbreak shapefile}. When false, the model derives the boundary from OpenStreetMap data using \texttt{city\_name}.
        \item \texttt{name\_shapefile}, the file name of the boundary shapefile, read from the \texttt{shapefile\_\allowbreak base} directory, used when \texttt{use\_\allowbreak shapefile} is true.
    \end{itemize}

    \item \emph{POI source selection}:
    \begin{itemize}
        \item \texttt{poi\_from\_shp}, a flag that tells the model which of the two ways to obtain \emph{POI} data to use. When true, the model reads \emph{POI} information from the shapefiles listed in \texttt{poi\_shapefile\_paths}. When false, the model downloads \emph{POI} data from OpenStreetMap.
        \item \texttt{poi\_shapefile\_paths}, the list of file paths to the \emph{POI}
        shapefiles, used when \texttt{poi\_from\_shp} is true.
        \item \texttt{poi\_label\_field}, the name of the attribute in the input \emph{POI}
        shapefiles whose values the model uses to associate each record with one or
        more configured \emph{POI types}.
    \end{itemize}

    \item \emph{radius filtering}:
    \begin{itemize}
        \item \texttt{poi\_radius\_enabled}, a flag. When true, the
        model pre-filters candidate POI access locations to those that lie
        within the \emph{global extraction radius} of at least one origin node,
        reducing the routing problem before any impedance computation is performed.
        Its default value is true.
        \item \texttt{poi\_radius\_decay\_threshold}, the minimum meaningful decay
        value used to derive the \emph{global extraction radius} from the configured
        decay coefficients. Its default value is \(0.5\).
        \item \texttt{poi\_radius\_max\_speed\_kmh}, the reference travel speed, in
        kilometres per hour, used to convert the threshold travel time to a distance
        when deriving the \emph{global extraction radius}. Its default value is
        \(20\).
    \end{itemize}

    \item \emph{impedance parameters}: Impedance is the model's measure of how costly a trip is, expressed as a single time value in minutes so that different trips can be compared on the same scale. These fields adjust how the model computes it:
    \begin{itemize}
        \item \texttt{drive\_access\_time\_min}, a fixed number of minutes the model adds only to car trips, to account for the overhead of accessing a vehicle. Walking, cycling, and public-transport trips do not receive this overhead. Its default value is \(10\).
        \item \texttt{cost\_per\_liter} and \texttt{distance\_for\_liter}, which together set the fuel cost of a car trip only, as a price per liter of fuel and the distance a liter covers. Walking, cycling, and public-transport trips do not receive this cost. Their default values are \(1.8\) and \(15\).
        \item \texttt{vot} (value of time), the coefficient the model uses to convert a monetary cost into an equivalent number of minutes. Its default value is \(0.2\).
    \end{itemize}


    \item \emph{transport scope}:
    \begin{itemize}
        \item \texttt{gtfs\_feeds}, the ordered list of file paths to the \emph{GTFS feeds}
        used for public-transport routing.
        \item \texttt{bus\_departure\_dt}, the date and time the model uses for public-transport
        routing.
        \item \texttt{bus\_ticket\_price}, \texttt{metro\_ticket\_price}, and
        \texttt{train\_ticket\_price}, the fare of a single journey on each public
        transport mode. Their default values are \(1.3\), \(2.55\), and \(2.55\).
        \item \texttt{time\_indifference\_bus}, a reference number of minutes the
        public-transport routing stage uses to judge how much a fare should discount
        the perceived weight of in vehicle time. Its default value is \(60\).
    \end{itemize}
    \notimpl{
    For each public-transport journey, the routing stage computes a discount factor,
    \(\gamma\), from \texttt{time\_\allowbreak indifference\_\allowbreak bus} (\(T\)), \texttt{vot} (\(v\)),
    and the relevant ticket price (\(p\)):
    \begin{equation}
        \gamma = \frac{T - v \cdot p}{T}
    \end{equation}
    The stage then applies \(\gamma\) only to the in vehicle time (\(t_{\text{ride}}\)),
    and adds the fare's time equivalent (\(v \cdot p\)) again on its own, alongside the
    other trip components:
    \begin{equation}
        \phi_{\text{bus}}(x,y) = t_{\text{access}} + t_{\text{wait}} + \gamma \cdot t_{\text{ride}} + t_{\text{transfer}} + t_{\text{egress}} + v \cdot p
    \end{equation}
    The net effect of the fare on \(\phi_{\text{bus}}\) is
    \begin{equation}
        v \cdot p \cdot \left(1 - t_{\text{ride}}/T\right).
    \end{equation}
    For a journey whose in vehicle time is shorter than \texttt{time\_\allowbreak indifference\_\allowbreak bus},
    the fare adds a small amount to the impedance, as expected. For a journey whose in
    vehicle time exceeds \texttt{time\_\allowbreak indifference\_\allowbreak bus}, the fare instead lowers the
    impedance, so a longer ride with a fare can end up with a smaller impedance than the
    same ride without one.
    }

    \item \emph{capability aggregation}: these are fixed constants used by every run, not fields an analyst sets per analysis:
    \begin{itemize}
    \item \texttt{ELECTRE\_Q\_FACTOR}, the factor used to derive the
    ELECTRE TRI indifference threshold from the standard deviation of the
    service-score vector of the current origin node and \emph{capability}.
    Its value is \(0.25\).
    \item \texttt{ELECTRE\_P\_FACTOR}, the factor used to derive the
    ELECTRE TRI preference threshold, the same way.
    Its value is \(0.75\).
    \item \texttt{ELECTRE\_LAMBDA\_CUT}, the minimum credibility score two
    categories must reach for a capability to be judged as belonging to the
    higher one, in the ELECTRE TRI outranking comparison. Its value is \(0.65\).
    \end{itemize}

    \item \emph{visualisation specification}: these fields control how the model automatically styles the generated QGIS project:
    \begin{itemize}
        \item \texttt{default\_capability}, the \emph{capability} or \emph{service}
        field selected when the QGIS project is first opened.
        \item \texttt{show\_basemap}, a flag that adds a background map layer
        (CartoDB Positron, a light basemap) to the QGIS project when true.
        \item \texttt{hexagon\_radius}, the inradius of each regular \emph{hexagon},
        that is, the distance from its centroid to the midpoint of one of its sides.
        The model uses this same value both to sample origin nodes into a hexagonal
        grid and to build the hexagonal grid layer in the QGIS project. Its default
        value is \(100\) meters.
        \begin{figure}[H]
            \centering
            \includegraphics[width=0.3\linewidth]{hex.png}
        \end{figure}
        \item \texttt{capability\_colors}, the list of signature colours the model
        assigns to the configured \emph{capabilities}, one colour per capability, in
        the same order the capabilities are defined. With the case study's three
        capabilities, the default list is blue (\texttt{\#006BFF}) for restorativeness,
        orange (\texttt{\#FFA200}) for nutrition, and pink (\texttt{\#EB4CCC}) for care.
    \end{itemize}
    Every \emph{service} under a \emph{capability}, and the capability's own grid
    layer, use that same colour. The model computes each shade by blending white
    with the full signature colour, using blend ratios of \(20\%\), \(40\%\),
    \(60\%\), \(80\%\), and \(100\%\), one ratio per ELECTRE category from lowest
    to highest.


\end{itemize}

\subsubsection*{POI types configuration}
Each \emph{POI type configuration} record defines one conceptual class of destination. Its fields are:
\begin{itemize}
    \item \texttt{poi\_type}, the unique identifier of the \emph{POI type}.
    \item \texttt{decay\_coefficient}, the impedance value, in minutes, at which access to the \emph{POI type} reaches a decay value of \(0.5\).
    \item \texttt{choquet\_interactions}, the ordered interaction coefficients between this \emph{POI type} and the other \emph{POI types} that belong to the same \emph{service}. The position that corresponds to the \emph{POI type} itself is left empty. The order of the list follows the order of the \emph{POI type} list declared by the associated \emph{service}.
    \item \texttt{tags}, the OpenStreetMap tags the model queries to retrieve the \emph{POI type}. The model uses this field only when \texttt{poi\_from\_shp} is false, that is, when \emph{POI} data comes from OpenStreetMap rather than from shapefiles.
    \item \texttt{labels}, a list of strings that count as this \emph{POI type} in the column named by \texttt{poi\_label\_field}. A record whose value in that column matches any of these strings is associated with the configured \emph{POI type}, which lets one \emph{POI type} group several source categories together.
\end{itemize}

\subsubsection*{Service configuration}
Each service configuration record defines an ordered aggregation of POI types. Its fields are:
    \begin{itemize}
        \item \texttt{service}, the unique identifier of the \emph{service}.
        \item \texttt{poi\_types}, the ordered list of \emph{POI type} identifiers that contribute to that \emph{service}.
        \item \texttt{choquet\_capacity}, an ordered list of capacity values, one for each \emph{POI type} in \texttt{poi\_types}, used by the Choquet aggregation described later.
        \item \texttt{contribution\_coefficient}, an ordered list of coefficients, one for each \emph{POI type} in \texttt{poi\_types}, each indicating how many accessible \emph{POIs} of that type are needed for the accessibility value to reach \(0.9\).
    \end{itemize}
A \emph{service configuration} record is only valid when \texttt{poi\_types}, \texttt{choquet\_\allowbreak capacity}, and \texttt{contribution\_\allowbreak coefficient} all have the same length, and every \emph{POI type} it references exists in the \emph{POI type configuration}.
\subsubsection*{Capability configuration}
    Each capability configuration record defines an ordered aggregation of
    services. An analyst can add, remove, or rename \emph{capabilities}
    entirely through this file. Its fields are:
    \begin{itemize}
        \item \texttt{capability}, the identifier of the \emph{capability}.
        \item \texttt{services}, the ordered list of \emph{services} that contribute to that capability.
        \item \texttt{electre\_weight}, an ordered list of weights, one for each
        \emph{service} in \texttt{services}, used by the ELECTRE TRI aggregation
        described later.
        \item \texttt{veto\_threshold}, the veto threshold used by the ELECTRE
        TRI aggregation for this capability.
        \item \texttt{enabled}, a flag. When false, the model skips computing
        this capability entirely: its score is left blank in the output, and no
        grid layer, service layer, or legend is generated for it in the QGIS
        project.
    \end{itemize}
    A \emph{capability configuration} record is only valid when every
    \emph{service} it references exists in the \emph{service configuration}.



\subsection{Intermediate objects}
\subsubsection*{Impedance bundle}
    The \emph{impedance bundle} is a cache artifact that combines the public
    transport impedance matrix and the walking, cycling, and driving routing
    caches into a single file (\texttt{impedances.npz}), so a later run can
    skip routing entirely when nothing relevant has changed. Its fields are:
    \begin{itemize}
        \item \texttt{schema\_version}, the format version of the bundle, used
        to detect a bundle produced by an older, incompatible version of the
        model.
        \item \texttt{node\_ids}, the identifiers of the origin nodes the
        bundle covers.
        \item \texttt{bus\_impedance\_matrix}, the dense origin by destination
        public-transport impedance matrix.
        \item \texttt{bus\_dest\_coords}, the coordinates of each destination
        access location in the bus impedance matrix.
        \item \texttt{routing\_departure\_iso}, the departure date and time
        the public-transport routing used.
        \item \texttt{origins\_sig} and \texttt{destinations\_sig}, stable
        signatures of the ordered origin and destination sets, used to detect
        whether the current run's node set still matches the one the bundle
        was built for.
        \item \texttt{subway\_impedance\_matrix} and
        \texttt{subway\_dest\_coords}, the same two fields for the subway
        modality, present only when subway routing is enabled.
        \item one walking, cycling, and driving routing cache per origin
        node, stored as an individual entry inside the bundle.
    \end{itemize}
    The model reuses a bundle only when its \texttt{schema\_version},
    \texttt{origins\_sig}, and subway coverage all match the current run;
    otherwise it treats the bundle as stale and recomputes it.

\subsubsection*{Global extraction radius}
    The \emph{global extraction radius} is a scalar distance in metres the
    model uses to pre-filter candidate POI access locations before any
    routing computation is performed. The model derives it from the
    configured \emph{POI type} decay coefficients: it selects the largest
    \texttt{decay\_coefficient} across all configured \emph{POI types},
    computes the travel time at which the corresponding decay function falls
    to \texttt{poi\_radius\_decay\_threshold}, and converts that time to a
    distance using \texttt{poi\_radius\_max\_speed\_kmh}.

\subsubsection*{Bus routing result}
    The \emph{bus routing result} is a dense matrix that gives one public
    transport impedance value, in minutes, for every origin node and
    destination access location. The model computes it in a single pass over
    the routing matrix produced by R5, the public-transport routing engine:
    for each origin-destination pair, it reads the trip's time breakdown and
    fare and computes the impedance using the formula described in \emph{transport
    scope}.

    The result is stored as:
    \begin{itemize}
        \item the impedance matrix itself,
        \item an index mapping origin identifiers to their row position in
        the matrix,
        \item an index mapping destination identifiers to their column
        position in the matrix.
    \end{itemize}
    A value of zero in the matrix means the destination is unreachable from
    that origin, either because no transit connection exists or because the
    only route found was a walk-only path.

\subsubsection*{Walking, cycling, and driving routing results}
    The \emph{walking, cycling, and driving routing result} is defined for one
    origin node. It groups its data by \emph{service}, and within each
    service, by \emph{POI type}. For each \emph{POI type}, it holds:
    \begin{itemize}
        \item the coordinates of every matching point of interest,
        \item the coordinates of the access location selected for each one,
        \item their source \emph{POI} identities,
        \item three equally aligned impedance vectors, for walking, cycling,
        and driving, each value expressed in minutes, and null when a
        candidate is unreachable by that mode,
        \item the walkability score aligned with each walking impedance.
    \end{itemize}
    The same point of interest can appear under more than one service, since
    its \emph{POI type} entry is duplicated under every service it belongs
    to.

\subsubsection*{Walkability scores for edges}
    The edge walkability object is a relation from a directed network-edge
    identifier to a walkability score on a scale from one to five. Each cached
    set of edge scores is tagged with a signature computed from the walking
    network it was built from. If the network is rebuilt and changes, that
    signature changes too, and the model discards the old scores instead of
    reusing ones that no longer match the current streets. The walkability of
    a path is obtained as a length-weighted aggregation of the scores of the
    edges that compose it.

    Each edge's score is meant to be the average of three separate scores,
    one for a point sampled at the start of the edge, one at its midpoint,
    and one at its end. Each of those three scores comes from an image
    classification model, a Swin Transformer, that predicts a walkability
    class from one to five for a street level image taken at that point.
    \notimpl{This scoring has not been implemented yet: every walking edge is
    currently given the same, highest possible score, so the walkability of
    every walking path currently comes out the same regardless of the streets
    actually used.}

\subsubsection*{Accessibility results}
    An accessibility result is defined for one origin node. Its fields are
    \texttt{node\_id}, \texttt{lat}, \texttt{lon}, \texttt{accessibility\_\allowbreak by\_\allowbreak service},
    a mapping from each \emph{service} to an ordered collection of \emph{POI type}
    accessibility records, and \texttt{accessibility\_\allowbreak by\_\allowbreak poi}, a mapping from each
    individual point of interest to its own accessibility value. Each \emph{POI type}
    accessibility record holds a \emph{POI type} identifier and an accessibility
    value in $[0,1]$.

    For the access location selected for each \emph{POI} and origin node, the modal
    impedance values, walking, cycling, driving, and public transport, are
    transformed into modal decay values through a decay function parameterised per
    \emph{POI type}, so that the \emph{POI type}'s configured reference impedance
    corresponds to a decay of $0.5$. The modal decays of a candidate are combined
    into a single multimodal route resource availability value.
    \texttt{accessibility\_by\_poi} keeps this value for each source \emph{POI};
    the \emph{POI type} accessibility record instead reduces the ordered collection
    of these values for a \emph{POI type} into a single \emph{POI type}
    accessibility value.

\subsubsection*{Service results}
    A service result is defined for one origin node. Its fields are
    \texttt{node\_id}, \texttt{lat}, \texttt{lon}, and \texttt{service\_scores},
    a mapping from every configured \emph{service} identifier to one normalised
    service score in $[0,1]$. Each score aggregates the accessibility values of
    the ordered \emph{POI types} declared by that \emph{service}, weighting each
    \emph{POI type} by its singleton capacity and accounting for the pairwise
    interactions between the \emph{POI types} of the service.


\subsection{Output objects}
The model produces two output objects: the \emph{capability results} and the \emph{visualization map}. The \emph{capability results} store the computed capability ratings and the underlying \emph{service} ratings in one tabular file, one record per point. The \emph{visualization map} associates each record with its point geometry and builds a hexagonal grid for map based visualization. Users explore both outputs through the \emph{provenance interface}, which adds interactive behaviour to the \emph{visualization map} and lets users inspect the provenance of each \emph{capability} or \emph{service} rating, including the \emph{POIs} used in the analysis and their individual contributions.
\subsubsection*{Capability results}
The capability results are stored in a single CSV file. Each row corresponds to a \emph{node} in the road graph and is defined by:
\begin{itemize}
    \item \texttt{node\_id}, which uniquely identifies the \emph{node}.
    \item \texttt{lat} and \texttt{lon}, the geographic coordinates of the \emph{node}.
    \item one \texttt{capability\_\textnormal{\emph{name}}} field for each enabled \emph{capability}, storing the rating of the represented \emph{node}. A disabled capability has no field at all.
    \item one \texttt{service\_\textnormal{\emph{name}}} field for each \emph{service} contributing to any \emph{capability}, storing the individual \emph{service} rating for the \emph{node}.
\end{itemize}
The \texttt{node\_id} links each record to the corresponding service-level ratings in the \emph{capability results} and to its POI-level information in the associated provenance data structures.
\subsubsection*{Visualization map}
The \emph{visualization map} is made of \emph{coloured hexagons} and their associated \emph{representative points}.
Each \emph{representative point} is defined by:
    \begin{itemize}
        \item \texttt{node\_id}, which uniquely identifies the \emph{node} represented by the point.
        \item \texttt{lat} and \texttt{lon}, the geographic coordinates of the represented \emph{node}.
        \item one \texttt{capability\_\textnormal{\emph{name}}} field for each enabled \emph{capability}, storing the rating of the represented \emph{node}.
    \end{itemize}
Each \emph{coloured hexagon} is defined by
\begin{itemize}
    \item \texttt{size}, which is determined by the \texttt{hexagon\_radius} in the \emph{general configuration}
    \item \texttt{centroid}, the geometric center of the \emph{hexagon}
    \item \texttt{representative\_point}, the \emph{node} in the road graph closest to the hexagon's \texttt{centroid}, used to represent all points that fall within that \emph{hexagon}.
    \item \texttt{colour}, which is determined by the rating associated with the \texttt{representative\_point}.
\end{itemize}
The \emph{representative points} are stored in the map but are hidden by default; only the \emph{coloured hexagons} are visible.
\begin{figure}[H]
    \centering
    \includegraphics[width=0.7\linewidth]{capability_viz.png}
    \caption{An example of a visualization map}
    \label{fig:vizmap}
\end{figure}

\subsubsection*{Provenance interface}
The \emph{provenance interface} allows interactive exploration of the \emph{visualization map}. Its behaviour is controlled by the following interface state fields:
\begin{itemize}
    \item \texttt{current\_mode}, which indicates whether the interface is in \emph{capability overview mode}, \emph{hexagon detail mode}, or \emph{service detail mode}.
    \item \texttt{mode\_description}, a textual description of the current mode, explaining what is shown and how the current visualization should be interpreted.
    \item \texttt{current\_capability}, the selected \emph{capability} whose ratings are shown in the \emph{visualization map}.
    \item \texttt{scale}, the conversion between real distance and distance in the visualization.
    \item \texttt{zoom\_level}, the current zoom level of the visualization.
    \item \texttt{heatmap\_legend}, based on the interval division of the capability ratings.
    \item \texttt{available\_services}, the \emph{services} that are connected to the \texttt{current\_capability} and are represented as selectable interface controls.
    \item \texttt{current\_service}, the selected \emph{service} whose ratings and POI contributions are shown when the interface is in \emph{service detail mode}.
    \item \texttt{hovered\_service}, the \emph{service} currently hovered by the user in \emph{hexagon detail mode} or \emph{service detail mode}. When this field is not null, the interface highlights the POIs that contribute to the hovered \emph{service}.
    \item \texttt{show\_colours}, whether or not the current colour visualization is active.
\end{itemize}
When the interface is in \emph{hexagon detail mode} or \emph{service detail mode}, it also shows POI information related to the selected \emph{hexagon} and \emph{capability}. In \emph{service detail mode}, this information is filtered according to the selected \emph{service}. Each \emph{POI} is defined by:
\begin{itemize}
    \item \texttt{poi\_id}, the unique identifier used to refer to the \emph{POI}.
    \item \texttt{poi\_coordinates}, the geographic coordinates of the \emph{POI}, expressed in angular form.
    \item \texttt{poi\_address}, the address of the \emph{POI}, expressed in terms of street, neighbourhood, and CAP.
    \item \texttt{type\_services}, the key-value pairs describing the conceptual POI types and services to which the \emph{POI} contributes. Each \emph{POI} can be associated with multiple conceptual POI types, which may contribute to different \emph{services}.
    \item \texttt{service\_power}, for each \emph{service} the \emph{POI} contributes to, the service power is computed as the accessibility value of the \emph{POI} multiplied by the capacity of the corresponding conceptual POI type. Before export, the model rescales these values within each \emph{service} by empirical rank, so that a distribution of powers that is tightly clustered in practice is spread evenly across $[0,1]$ instead of looking nearly uniform.
    \item \texttt{capability\_power}, for each \emph{capability} to which the
    \emph{POI} contributes, an explanatory value derived from the POI's
    contributions to the relevant \emph{services} and from the service weights
    used in the ELECTRE TRI aggregation, rescaled the same way, within each \emph{capability}, before export.
   \item \texttt{is\_selected}, which indicates whether the \emph{POI} contributes to the \texttt{current\_service} when the interface is in \emph{service detail mode}, or to the \texttt{hovered\_service} when a service is hovered in \emph{hexagon detail mode}. If selected, the point is highlighted with a colour; otherwise, it is shown with transparency.
    \item \texttt{point\_size}, the size assigned to the visual representation of the \emph{POI}. It is mostly fixed, with a modest increase when the \emph{POI} is hovered, highlighted through a hovered \emph{service}, or coloured by \emph{POI type}; it does not vary continuously with power.
    \item \texttt{point\_colour}, which is assigned to the \emph{POI} when \texttt{show\_colours} is active, as a greyscale gradient from light, for a low power value, to dark, for a high one. In \emph{hexagon detail mode} the shade reflects \texttt{capability\_power}; in \emph{service detail mode} it reflects \texttt{service\_power}. Hovering a \emph{service} overrides this: POIs that contribute to the hovered \emph{service} are highlighted in green, and the rest are shown with reduced opacity.
\end{itemize}


\section{Dynamics}
\subsection{Input processing}
When a \emph{boundary} is provided, the model validates it, normalises it to the spatial reference system used by the analysis, and derives a single study area geometry from it. This geometry defines the spatial extent for graph extraction and for all subsequent origin-destination computations. When no explicit \emph{boundary} is provided, the model derives the study area from the configured \texttt{city\_name}. The model loads the available public-transport feeds and identifies the services that are active for the configured routing date and time. When multiple feeds are provided, they are considered together so that public-transport impedance can account for the available transport modes.

\subsection{Network preparation}
A separate routable network is prepared for each travel mode considered by the model. Each network is spatially limited to the selected study area and preserves the nodes and edges needed to compute origin-destination impedance. The resulting mode-specific networks provide the routing substrate for the subsequent impedance calculation.

The model evaluates capability at every node of the walking network, since walking is the one mode available to any traveler regardless of what other transport they can use, and treats those nodes as the origins for the whole analysis. To keep this origin set at a manageable, spatially even density rather than the raw number of street intersections, the model groups these nodes into a regular hexagonal grid and keeps only one representative node, the one closest to each hexagon's centre, per occupied cell. A larger \texttt{hexagon\_radius} covers more ground per hexagon, so fewer, more widely spaced origins are kept; a smaller \texttt{hexagon\_radius} keeps origins closer together, at the cost of evaluating more of them.

\subsection{POI preparation}
When POIs are provided through external files, the model loads the available point, line, and polygon records and associates each record with one or more configured \emph{POI types} using the record label. Records that cannot be associated with any configured \emph{POI type} are excluded from the analysis. When POIs are downloaded from OpenStreetMap, the model retrieves the POIs required by the configured \emph{POI types} within the study area and assigns them to the corresponding conceptual classes.
\subsection{POI snapping}
Each POI geometry is converted into one or more candidate access locations. A point geometry defines a single access location, while a line or polygon geometry defines multiple candidate access locations, one for each vertex composing it: every vertex of the line, or every vertex of the polygon's boundary. Each candidate access location is associated with the nearest node of the relevant travel mode network. The resulting snapped nodes represent the network entry points used during routing. When a \emph{POI} has multiple candidate access locations, the model selects one candidate separately for each source node. The selected candidate is the one with the smallest Haversine distance from the source. This allows a spatially extended \emph{POI} to be represented by different access locations depending on the source from which it is evaluated. When \texttt{poi\_radius\_enabled} is enabled, candidate access locations
that do not lie within \(r_{\text{global}}\) of any origin node are discarded
before the origin specific candidate selection is performed.
\subsection{Extraction radius}
When \texttt{poi\_radius\_enabled} is set, the model derives a single
\emph{global extraction radius} $r_{\text{global}}$ before any routing is
performed. The radius is computed from the \emph{POI type} $k$ with the largest
configured \texttt{decay\_coefficient} $d_k$, which gives the most permissive bound
and conservatively covers every type. Given that decay coefficient and the
configured \texttt{poi\_radius\_decay\_threshold} $\tau$, the maximum travel
time $t_{\max}$ beyond which decay falls below $\tau$ is

\begin{equation}
    t_{\max} = \frac{-\ln \tau}{\beta_k}, \qquad \beta_k = \frac{\ln 2}{d_k}
\end{equation}

\noindent and the corresponding straight line distance is

\begin{equation}
    r_{\text{global}} = t_{\max} \cdot \frac{v_{\max}}{60} \cdot 1000
\end{equation}

\noindent where $v_{\max}$ is \texttt{poi\_radius\_max\_speed\_kmh}.
When \texttt{poi\_radius\_enabled} is enabled, the \emph{global extraction
radius} is used to pre-filter candidate POI access locations during
\emph{POI snapping}, \emph{public-transport routing}, and
\emph{walking, cycling and driving routing}. When radius filtering is
disabled, no distance based pre-filter is applied.
\subsection{Routing}

\subsubsection*{Public transport routing}
Unlike walking, cycling, and driving, public transport routing does not use a
dedicated network of its own: it relies on the walking network, and the
access locations snapped during \emph{POI snapping} for the walking mode are
reused as its destination candidates. When \texttt{poi\_radius\_enabled} is enabled, the full destination set is
reduced before routing to the access locations that fall within the
\emph{global extraction radius} of at least one origin node. When radius
filtering is disabled, the full set of candidate access locations is retained. The model then resolves the best available public-transport connection departing at the configured date
and time for every origin-destination pair, and reduces each connection to
a single impedance value using the formula described in \emph{transport
scope}. The resulting impedance matrix is
cached and reused in subsequent runs unless the origin set, the destination
set, or the departure time has changed.

\subsubsection*{Walking, cycling and driving routing}
For each origin node, the model determines independently for each travel mode
the shortest-path to the candidate access location selected for each
\emph{POI} with respect to the current origin node. When
\texttt{poi\_radius\_enabled} is enabled, only candidates whose straight line
distance from the origin does not exceed the \emph{global extraction radius}
are considered. When radius filtering is disabled, all relevant candidate
access locations are considered. Each of these three modes routes on its own dedicated network so that modal constraints such as permitted roads and design speeds are respected. Public transport is the exception: rather than routing on a dedicated network of its own, it uses the walking network, since pedestrian access to and from transit stops is assumed to follow the same streets a walking trip would use. Access locations that are not
reachable from the origin in a given modal network receive no routing result
for that mode. Results are stored per origin node and reused in subsequent
runs.

\subsection{Impedance}
The model derives a scalar impedance value from each routing result, expressing
the cost of reaching a destination as a single quantity in minutes that is
comparable across modes and usable in the decay function.

For public transport, the model computes the impedance directly from the
routing result using the formula described in \emph{transport scope}.

For walking, cycling, and driving, the impedance combines a quality
adjustment with a time and cost component. The quality adjustment applies
only to walking. For a path with walkability score \(w \in [1,5]\), the
quality factor is the multiplicative coefficient

\begin{equation}
    \eta(w) = 1 + \lambda\,\frac{5-w}{4},
    \qquad \lambda = 0.15.
\end{equation}

\noindent The walking impedance is therefore

\begin{equation}
    \phi_{\text{walk}}(x,y)
    = \eta(w)\,t_{\text{walk}}(x,y),
\end{equation}

\noindent where \(t_{\text{walk}}(x,y)\) is the shortest-path travel time
between origin \(x\) and destination \(y\) on the walking network. The quality
factor equals \(1\) when walkability is maximal (\(w=5\)) and \(1.15\) when
walkability is minimal (\(w=1\)). When the walkability score is unavailable,
\(\eta(w)\) defaults to \(1\). \notimpl{Since edge scores are not yet
implemented (see \emph{Walkability scores for edges}), every walking path is
currently treated as maximally walkable, so this adjustment has no visible
effect on results yet: \(\eta(w) = 1\) for every walking connection.}

For cycling, the impedance is simply the shortest-path travel time on the
cycling network, with no additional overhead or cost.

For driving, the impedance adds a fixed access overhead
(\texttt{drive\_access\_time\_min}) and an equivalent time cost for fuel to
the shortest-path travel time:

\begin{equation}
    \phi_{\text{drive}}(x,y)
    =
    \texttt{drive\_access\_time\_min}
    +
    t_{\text{drive}}(x,y)
    +
    \texttt{vot} \cdot \frac{\texttt{cost\_per\_liter}}{\texttt{distance\_for\_liter}}.
\end{equation}

Connections for which no transit service exists, or for which only a walk-only
path was found, receive no impedance and do not contribute to accessibility.
Access locations that are disconnected from the origin in a given modal graph
also receive no impedance for that mode.

\subsection{Decay}
The model converts each impedance value into a decay value through an
exponential function. For a \emph{POI type} \(k\) with decay coefficient \(d_k\), the decay rate is

\begin{equation}
    \beta_k = \frac{\ln 2}{d_k},
\end{equation}

\noindent so that the decay reaches \(0.5\) at the configured reference
impedance \(d_k\).
Given an impedance \(\phi^i_{m,k}\), the corresponding
decay is
\begin{equation}
    \psi_{m,k}(\phi_{m,k}^i)(x,y) = e^{-\beta \phi^i_{m,k}}.
\end{equation}

When a routing result is unavailable, the corresponding decay is treated as
zero. Modal decay values are then combined into a single route resource
availability value, so that a destination is not penalised simply because one
mode is weak if another realistic mode provides access. The decay stage
therefore turns raw travel cost into a bounded accessibility signal that can be
used consistently across all later aggregation steps.

\subsection{Accessibility}
Accessibility is computed for each origin node and each configured
\emph{POI type}. The model considers the access location selected for each \emph{POI}
associated with that \emph{POI type} and the current origin node, retrieves the corresponding impedance values for walking, cycling, driving, and public transport, and converts those impedances
into modal decays with the \emph{decay} function. The modal decays of one
candidate access location are merged into a single value, and the candidate
values are then reduced to one accessibility value in the interval \([0,1]\).
The reduction is controlled by the \texttt{contribution\_coefficient}
associated with the \emph{POI type}. A larger coefficient means that more accessible POIs are required for the accessibility value to approach saturation, causing the score to grow more
gradually. Conversely, a smaller coefficient causes the accessibility value
to approach saturation with fewer accessible POIs. The resulting accessibility value therefore reflects both the
quality of the routes and the multiplicity of accessible destinations.

\subsection{Service aggregation}
The accessibility values of the \emph{POI types} belonging to the same
\emph{service} are aggregated with a Choquet integral. The input vector is first
aligned to the ordered \emph{POI type} list declared by the service
configuration, and the values are then accumulated from the smallest to the
largest accessibility. The aggregation uses the singleton capacities and the
pairwise interactions defined for that service, so that complementary
\emph{POI types} can reinforce one another while redundant combinations can
contribute less than a simple sum would suggest.
The service score is normalised by the measure of the full set of
\emph{POI types}, so every service score remains in \([0,1]\). This keeps the
service layer directly comparable across services and ensures that the
capability layer receives inputs on a common scale.

\subsection{Capability aggregation}
The service scores of each \emph{capability} are aggregated with the
ELECTRE TRI method. The ordered service score vector is compared against the
boundary profiles \(0.2\), \(0.4\), \(0.6\), and \(0.8\), using the service
weights declared in the \emph{capability configuration}. For each origin node
and \emph{capability}, the indifference threshold \(q\) and preference
threshold \(p\) are dynamically derived from the standard deviation of the
current service score vector \(x\):

\begin{equation}
    q = \operatorname{std}(x) \cdot \texttt{ELECTRE\_Q\_FACTOR},
    \qquad
    p = \operatorname{std}(x) \cdot \texttt{ELECTRE\_P\_FACTOR},
\end{equation}

\noindent The resulting \(q\) and \(p\) values are applied uniformly to all
\emph{services} belonging to the current \emph{capability}.

For each boundary, the model computes a weighted concordance score from how
far each service score sits above or below the boundary relative to \(q\)
and \(p\). It then checks whether any single service falls far enough below
the boundary to exceed the capability's \texttt{veto\_threshold}; if one
does, the concordance for that boundary is forced to zero regardless of how
the other services scored. The capability is assigned to the highest
boundary whose resulting credibility reaches \texttt{ELECTRE\_LAMBDA\_CUT}.

The result of the ELECTRE TRI procedure is a category assignment: this
five-way rating is the capability output. For storage and display it is
encoded as the midpoint of the assigned category interval, giving one of
five fixed values in \((0,1)\).

\subsection{Capability results visualization}
The final capability scores are written to a single CSV file and transformed into a
spatial \emph{visualization map}. Each retained representative node becomes a \emph{representative point}
carrying the capability and service scores computed for that node. The points are then used to build a regular hexagonal grid in projected coordinates, and each hexagon is coloured according to the
selected capability score of its representative point.
This representation provides both a node-level and an area-level view of
the results. The node points preserve the exact network based measurement
locations, while the hexagonal aggregation makes the spatial distribution of
the scores easier to inspect. The generated outputs are also used by the
\emph{provenance interface}, which allows the user to switch between
capabilities and inspect the POIs that contributed to each score.
\subsection{Provenance interface dynamics}
The \emph{provenance interface} supports the following general interaction actions:
\begin{itemize}
    \item The selected \emph{capability} can be changed through a dropdown menu, which updates the scores and colours shown in the \emph{visualization map}.
    \item A heatmap legend explains how to interpret the colours of the current visualization and is updated when the selected \emph{capability} requires a different interval division.
    \item The user can zoom in and out using the buttons located in the top left or by scrolling with the mouse.
    \item The interface allows the user to zoom out until all displayed points are visible and to zoom in until the shown area is contained within a single \emph{hexagon}.
    \item When zooming, the interface updates the scale visualizer in the bottom right.
    \item The mode description in the bottom left explains the current interaction mode and how the displayed information should be interpreted.
    \item On the right, a toggle allows the user to turn the colour visualization on or off, so the base map underneath can be seen.
\end{itemize}
The \emph{provenance interface} is organized into three interaction modes:
\begin{enumerate}
    \item \emph{Capability overview mode}. Initially, the interface is set to \emph{capability overview mode}. In this mode, the \emph{hexagons} are coloured based on the scores of the \texttt{current\_capability}. The map provides an overview of how the selected \emph{capability} varies across the study area. Selecting a \emph{hexagon} brings the interface to \emph{hexagon detail mode}.
    \item \emph{Hexagon detail mode}. In \emph{hexagon detail mode}, the interface zooms into the selected \emph{hexagon}. The POIs used to compute the capability score of the selected \emph{hexagon} are shown on the map, each shaded on a greyscale gradient by its \texttt{capability\_power}. The user also sees all the \emph{services} associated with the \texttt{current\_capability}, each represented as a button on the right side of the interface. Hovering the mouse over one of these buttons highlights the POIs that contribute to the corresponding \emph{service} in green and dims the rest. Clicking a \emph{POI} shows its angular geographic coordinates, street, neighbourhood, CAP, and the \emph{POI types} and \emph{services} to which it contributes. When a \emph{service} is hovered, the displayed information is filtered according to the \texttt{hovered\_service}. In both \emph{hexagon detail mode} and \emph{service detail mode}, a back button allows the user to return to \emph{capability overview mode}. Clicking a \emph{service} brings the interface to \emph{service detail mode}. The service buttons remain available in \emph{service detail mode}, allowing the user to change the selected \emph{service}.
    \item \emph{Service detail mode}. In \emph{service detail mode}, only the POIs that contribute to the selected \emph{service} remain shown, shaded on the same greyscale gradient, now by \texttt{service\_power} rather than \texttt{capability\_power}. Clicking a \emph{POI} shows the coordinates, street, neighbourhood, CAP, and the POI type associated with the selected \emph{service}.
\end{enumerate}

\section{Implementation}
\subsection{Input processing}
The boundary file is read from the configured path and its CRS is inspected. If the CRS differs from WGS84 (EPSG:4326), the geometry is reprojected. Multiple polygon features are dissolved into a single geometry by applying a spatial union. Invalid geometries, such as rings with self-intersections, are repaired using a zero-distance buffer operation. When no boundary file is provided, the configured \texttt{city\_name} is resolved through the OpenStreetMap geocoder.
Each configured GTFS feed is checked for existence and staged as a \texttt{.zip} archive into the routing data directory used by R5. The model itself does not parse the GTFS schema or filter service patterns by date; R5 does that internally, driven by the configured \texttt{bus\_departure\_dt}, which the model passes to it as an environment variable.

\subsection{Network preparation}
The mode-specific graphs are extracted with OSMnx from OpenStreetMap data. When a boundary is provided, the dissolved study area polygon is used as the extraction geometry directly. When no boundary is provided, the model geocodes \texttt{city\_name} and expands the resulting boundary by the \emph{global extraction radius} before extraction, so that POIs just outside the city's official limits remain reachable. Graph simplification is disabled so that all network nodes are retained. Each graph is serialized to disk in GraphML format after its first construction and loaded from cache on subsequent runs.

For routing, each mode's network is additionally converted into a lightweight lookup table: a compressed sparse row representation of the network's connectivity and edge lengths, together with the node coordinates needed to snap points onto it and, for walking, a walkability score aligned to each edge. This lightweight version is rebuilt automatically whenever the underlying network file changes. It uses a small fraction of the memory that the full graph representation would need, which is what makes it practical for every parallel worker process in later stages to keep its own copy of the whole network in memory at once.

\subsection{POI preparation}
When the POI input is a shapefile collection, all point, line, and polygon feature files are read and concatenated into a single feature set. Each file's CRS is checked and reprojected to WGS84 when necessary. For each configured \emph{POI type}, the dataset is filtered by comparing the values of the attribute identified by \texttt{poi\_label\_field} with the values listed in \texttt{labels}. Unmatched records are logged and excluded. When the POI source is OpenStreetMap, a single batched query is issued for the study area. The resulting dataset is cached as a GeoJSON file and filtered in memory by applying the tag clauses defined in the model configuration. Individual POI type results are also cached to avoid repeated filtering.

\subsection{POI snapping}
Line geometries are decomposed into their vertices. Polygon geometries are decomposed into the vertices of their exterior ring and any interior rings. Multi-part geometries are processed one constituent geometry at a time. Candidate coordinates are snapped to the nearest graph nodes through a vectorised batch operation, which also returns the Haversine distance between the original coordinate and the snapped node. Snap results are persisted in a per-type binary cache identified by a hash of the POI type key and travel mode. Empty snap payloads are treated as stale and trigger re-extraction. For transit routing, the bus stage reuses the walk-mode snap results rather than computing a separate snap pass, because pedestrian access to transit stops is assumed to follow the walking network. For each source node and each non-point \emph{POI}, the Haversine distance between the source node and every snapped candidate access location is
computed, and the candidate with the smallest distance is selected.

\subsection{Extraction radius}
When \texttt{poi\_radius\_enabled} is \emph{True}, the model calls
\texttt{get\_global\_radius\_m} once before the routing stages begin. If
\texttt{poi\_radius\_m} is set, that value is returned directly. Otherwise,
the function identifies the maximum \texttt{decay\_coefficient} \(d_k\)
across all configured \emph{POI types} and calls
\texttt{threshold\_radius\_m}, which computes

\begin{equation}
    \beta_k = \frac{\ln 2}{d_k},
    \qquad
    t_{\max} = \frac{-\ln \tau}{\beta_k},
\end{equation}

\noindent where \(\tau\) is
\texttt{poi\_radius\_decay\_threshold}. It then returns

\begin{equation}
    r_{\text{global}}
    =
    t_{\max}\cdot\frac{v_{\max}}{60}\cdot 1000.
\end{equation}
\subsection{Public transport routing}
Public transport routing is delegated to r5r, a Java based library
for realistic transit routing. The model invokes r5r through an
external R script launched as a subprocess via Rscript. The routing
data bundle required by r5r consists of one OpenStreetMap PBF extract
and one or more GTFS feeds placed in a common directory. When no PBF file is
present, the model auto-builds one by downloading the unsimplified road network
for \texttt{city\_name} with osmnx and converting the resulting OSM
XML to PBF using osmium.

Origins and destinations are serialised as CSV files with columns
\texttt{id}, \texttt{lon}, and \texttt{lat}. Before writing these files, when \texttt{poi\_radius\_enabled} is enabled,
the model applies a destination pre-filter: a candidate snapped coordinate is
retained only if its Haversine distance to at least one origin node does not
exceed \(r_{\text{global}}\). When radius filtering is disabled, all candidate
snapped coordinates are retained.
The R script resolves the expanded travel time
matrix and writes one CSV row per origin-destination connection, containing
\texttt{from\_id}, \texttt{to\_id}, \texttt{departure\_time},
\texttt{access\_time}, \texttt{wait\_time}, \texttt{ride\_time},
\texttt{transfer\_time}, \texttt{egress\_time}, \texttt{routes},
\texttt{n\_rides}, and \texttt{total\_time}. Connections whose route sequence
indicates a walk-only trip are treated as absent.

For each valid row, the model computes the discount factor \(\gamma\) from
\texttt{time\_indifference\_bus} (\(T\)), \texttt{vot} (\(v\)), and the
relevant ticket price (\(p\)), and the impedance, exactly as described in
\emph{transport scope}:

\begin{equation}
    \gamma = \frac{T - v \cdot p}{T}
\end{equation}

\begin{equation}
    \phi_{\text{bus}}(x,y) = t_{\text{access}} + t_{\text{wait}} + \gamma \cdot t_{\text{ride}} + t_{\text{transfer}} + t_{\text{egress}} + v \cdot p
\end{equation}

\noindent where \(t_{\text{access}}\), \(t_{\text{wait}}\), \(t_{\text{ride}}\),
\(t_{\text{transfer}}\), and \(t_{\text{egress}}\) are read directly from the
\texttt{access\_time}, \texttt{wait\_time}, \texttt{ride\_time},
\texttt{transfer\_time}, and \texttt{egress\_time} columns of the R script's
output.

The resulting values are written into a memory mapped matrix of shape
\(N_{\text{origins}} \times N_{\text{destinations}}\) indexed by integer row
and column identifiers derived from the input CSVs.

\subsection{Walking, cycling and driving routing}
Walking, cycling, and driving routing is computed independently for each origin node, as a shortest-path search over the
lightweight per-mode lookup table described in \emph{Network preparation}
rather than over the full network object. Each parallel worker keeps its own
copy of this lookup table in memory instead of sharing the full network with
the other workers: once routing begins, a worker touches large parts of the
network regardless of the origin it is processing, so sharing a single full
network object across many workers would still leave each one holding its
own effective copy, multiplying memory use across all workers combined.
Every origin's starting point on each mode's network
is looked up once, before routing begins, so this lookup is not repeated for
every worker. For each origin node, when \texttt{poi\_radius\_enabled} is enabled, the model
first applies the Haversine pre-filter to limit candidate POI access locations
to those within \(r_{\text{global}}\) straight-line distance. When radius
filtering is disabled, this pre-filter is skipped. The shortest-path
search is also stopped once it has explored somewhat beyond the POI radius,
so it does not keep exploring the network indefinitely past the destinations
that could actually matter.
The shortest-path search then determines the travel time from the origin to the reachable candidate access location selected for each \emph{POI}.
Travel time on each edge is computed from the edge length and the speed assigned to the corresponding travel mode. For the driving mode, the configured
\texttt{drive\_access\_time\_min} is added to the shortest-path travel time to
account for vehicle access overhead, and an equivalent time cost for fuel is
added, using the same \(\phi_{\text{drive}}\) formula described in
\emph{Impedance}:

\begin{equation}
    \phi_{\text{drive}}(x,y)
    =
    \texttt{drive\_access\_time\_min}
    +
    t_{\text{drive}}(x,y)
    +
    \texttt{vot} \cdot \frac{\texttt{cost\_per\_liter}}{\texttt{distance\_for\_liter}}.
\end{equation}

\noindent For cycling and walking,
neither an access time nor an economic cost is added. Access locations that are disconnected from the origin in a given modal graph receive no impedance value for that mode.

For the walking mode the model additionally queries a per-edge walkability
index that assigns a score in \([1,5]\) to each directed edge. The walkability
score of a path is the length-weighted mean of the scores of its constituent
edges. \notimpl{This is not yet implemented for individual paths: the same
score is currently used for every edge, so the model does not yet reconstruct
which streets each walking route actually used in order to average their
scores.}

Results for each origin node are persisted to a dedicated pickle file
identified by the node identifier. A cache entry is considered valid when its
\texttt{schema\_version} field matches the configured
\texttt{non\_bus\_\allowbreak cache\_\allowbreak schema\_\allowbreak version}. A record of the POI
configuration is also stored alongside each cache entry, but is not currently
checked before a cached entry is reused. Valid entries are not recomputed
across runs.

The number of parallel workers used for this stage is chosen to keep the
pipeline's total memory use within a set limit. A baseline number of workers
is worked out from how much memory is available divided by how much memory
one worker is expected to need, and this stage additionally applies its own,
lower cap on top of that baseline. This extra cap exists because some origins
have unusually many nearby points of interest to process, and processing one
of these busy origins can briefly use a large amount of memory in that
worker; keeping the number of simultaneous workers lower limits how many of
these memory spikes can happen at the same time. Each
worker is also restarted after processing a bounded number of origins, and
periodically frees up any memory it is no longer using between origins, so
that memory used by past, already-finished origins does not keep
accumulating over the course of a run.

\subsection{Accessibility}
The accessibility stage reconstructs one node-level accessibility result for
each origin node by combining the cached walking, cycling, and driving routing result with the bus
impedance matrix. The stage first opens a dense memory mapped accessibility
matrix when a compatible cache already exists; otherwise it creates a fresh
matrix initialised with \texttt{NaN} and stores the expected row and column
mappings alongside it. Only rows that are still incomplete are scheduled for
recomputation.

Each worker loads the cached walking, cycling, and driving routing result for one origin node and retrieves
the relevant bus impedance values using the origin and destination index maps.
The worker then computes the \emph{decay} of every available modal impedance,
merges the modal decays into route resource availability, and reduces the
result to one accessibility value per \emph{POI type}. When the same
\emph{POI type} appears in several services, the stage can deduplicate the
computation and reuse the same accessibility value in all affected services.

The stage writes the computed accessibility values back into the dense matrix
so that later runs can reuse completed rows without repeating the routing
logic. This makes the accessibility stage resumable and avoids recomputing
origin rows that are already present in cache.

\subsection{Service aggregation}
The service stage loads the accessibility payload for each origin node and
groups the \emph{POI type} accessibility values by service. When a compatible
service matrix cache exists, the stage reconstructs the node-level service
results directly from the memory mapped matrix; otherwise it computes the
service scores from scratch and writes a new cache.

For each service, the ordered accessibility vector is passed to the Choquet
aggregation helper, which sorts the values from smallest to largest,
accumulates them weighted by the fuzzy measure of the \emph{POI types} not
yet accumulated (built from each type's singleton capacity plus its pairwise
interactions with the rest, as declared in the \emph{service configuration}),
and normalises the result by the measure of the full \emph{POI type} set so
the final score stays in $[0,1]$. The resulting service scores are stored per
node together with the node's coordinates.

\subsection{Capability aggregation}
The capability stage reads the service scores and writes a single CSV file
containing the results of all enabled \emph{capabilities}. Each row corresponds
to one origin node and contains the node identifier, the node coordinates, one
capability score field for each enabled \emph{capability}, and the service score
fields required by those capabilities. The stage computes the capability scores
from their respective ordered service vectors, using the service ordering and
weights declared in the \emph{capability configuration}. This is a rating,
stored as the ELECTRE TRI band midpoint.

After the CSV file is written, it is moved into the \texttt{experiments}
folder with a name derived from the study's city name. The stage also
appends the average capability values to a recap CSV so that multiple runs can
be compared without re-reading the complete capability results file.

\subsection{Capability results visualization}
The visualization step consumes the capability results CSV file and turns it
into geospatial outputs. The pipeline builds a shapefile and a GeoPackage from
the CSV records. The GeoPackage contains a point layer for the evaluated
representative nodes, hidden by default, and, when enabled, a hexagonal grid
layer in which each occupied cell is associated with the capability values of
its representative node.

The grid is built in projected coordinates so that the cell size is measured in
metres. Once the spatial files have been written, the pipeline optionally builds
a QGIS project, applies the requested styling parameters, and opens the project
automatically when QGIS is available. This provides a direct path from the
computed capability scores to an inspectable cartographic view.

\subsection{Provenance data export}
The interface itself is generated from the QGIS project using the qgis2web
plugin, which produces a standalone OpenLayers web map. The per-POI records
shown by the \emph{provenance interface} (the identity, service power, and
capability power of each \emph{POI} contributing to a hexagon) are exported
separately from the GeoPackage/QGIS outputs, as a set of small, compressed
files for the standalone offline interface. Hexagons are grouped into shard
files by a pure function of the hexagon's own identifier (a fixed block of
neighbouring hexagons per shard), so the interface computes which shard file
a given hexagon is stored in directly from its identifier. Within a shard,
repeated text such as service and capability names is replaced by short
numeric codes explained once in a small manifest file, and numeric values are
rounded to a fixed precision, which keeps the exported files small. Each
shard file is written incrementally as its hexagons are produced, so
producing this export uses a bounded amount of memory regardless of how many
hexagons exist, and the finished shard payloads are compressed so the
interface loads and decompresses them directly in a web browser.
\end{document}